\subsection{The move-first recording input}
\label{sv:recorder-import}
For an instruction $r=(q,a)$ with data $(q_r,b_r,d_r)$, the recorder
of \cite[Lemma~2.1]{Entry} is written here with its frontier symbol
$P$ renamed $F$ and its placement controls $F_q$ renamed $P_q$;
all other controls and guards are unchanged. It uses letters $(a,h,m)$,
$h\in\{E,F\}\sqcup\{[r]\}$ and $m\in\{0,1\}$,
with modes $S_q,A_r,R_r,P_q,L_q$.  Here is its complete table;
every occurrence of $h$ requires $h\ne F$:
\[
\begin{array}{c|c|c|c|c}
S_q&(a,h,0)&(b_r,h,0)&d_r&A_r\\
A_r&(c,h,0)&(c,h,1)&1&R_r\\
R_r&(c,h,0)&(c,h,0)&1&R_r\\
R_r&(c,F,0)&(c,[r],0)&1&P_{q_r}\\
P_q&(c,E,0)&(c,F,0)&-1&L_q\\
L_q&(c,h,0)&(c,h,0)&-1&L_q\\
L_q&(c,h,1)&(c,h,0)&0&S_q.
\end{array}
\]
There are no first-row instructions for halting $q$.  The cited lemma proves
that the target mode determines the incoming move, and the target mode
with the written complete letter determines the unique predecessor.
Initialize at the original machine head with frontier at any fixed
$r_0\ge2$, empty history elsewhere, and zero marks everywhere.
After $n$ machine steps the frontier is $r_0+n$; the work state,
tape and head are recovered at the $S$-visits.  If the next machine
head is $h'$, its next checkpoint takes exactly
$2(r_0+n-h')+4$ local transitions.  In particular these counts
are finite and positive.  Halting and indefinitely defined nonhalting
runs are preserved.  We use $r_0=2$ below.

Section~5.1 of~\cite{Entry} proves two extensions on the whole partial
domain: toggle the current mark only in ready states to obtain the
marker-preserving version, or copy an arbitrary passive track at every
write.  The latter permits a permanent origin flag.  We specify the
two changed rows again when the whole-box application uses the former.
Every other recorder in this paper has its own complete inverse proof.
